\documentclass[aps,pre,reprint,10pt,superscriptaddress]{revtex4-2}

\usepackage[T1]{fontenc}
\usepackage{lmodern}

\usepackage{mathtools,amssymb}

\usepackage{graphicx}
\usepackage{xcolor}
\usepackage{tikz}
\usetikzlibrary{calc, positioning}

\usepackage{booktabs}
\usepackage[caption=false]{subfig}

\usepackage{silence}
\WarningFilter{nameref}{The definition of \label has changed!}

\usepackage[hidelinks,pdfusetitle]{hyperref}
\usepackage[all]{hypcap}
\usepackage{cleveref}

\begin{document}

\title{Evolutionary adaptation through bet-hedging in finite populations under fluctuating environments}
\author{Marco Mendívil Carboni}
\author{Luis Dinis}
\author{Juan José Mazo}
\affiliation{Grupo Interdisciplinar de Sistemas Complejos (GISC) and Departamento de Estructura de la Materia, Física Térmica y Electrónica, Universidad Complutense de Madrid, 28040 Madrid, Spain}

\begin{abstract}
    Natural populations evolve under fluctuating environments and limited resources, yet the combined effects of these factors on adaptation remain unclear. Here, we introduce a minimal stochastic model of a population undergoing a birth--death process in a randomly switching environment, where individuals inherit phenotypic \emph{strategies} that determine the distribution of phenotypes across generations. We first confirm that environmental variability can favor phenotypic diversity (bet-hedging), and that the strategy maximizing the average growth rate generally differs from the one minimizing the extinction rate. We then show that evolution drives populations to converge to strategies that balance these competing objectives. We further demonstrate that population size modulates the evolutionary outcome: large populations recover the growth-maximizing strategy, whereas smaller ones evolve toward more resilient strategies. Finally, we show how a simple heuristic approximation captures these outcomes and clarifies the interplay between growth, extinction and mutation.
\end{abstract}

\maketitle

\section{Introduction and objectives} \label{sec:introduction}

Adaptation in fluctuating environments is a central problem in evolutionary biology \cite{gillespieNaturalSelectionVariances1977,kingscollegeRiskUncertaintyEvolutionary1982,kussellPhenotypicDiversityPopulation2005,ashcroftFixationFinitePopulations2014,saakianAlleleFixationProbability2019,camachomateuPhenotypicdependentVariabilityEmergence2021,kaushikTimeFixationChanging2021,collectiveEconomicPrinciplesCell,nakashimaPopulationDynamicsModels2023}. Natural populations rarely experience constant conditions: rather, they face environments that change \emph{unpredictably} over time, often on evolutionary timescales. In such contexts, if sensing the changes in the environment is infeasible or too costly, it is well established in the literature that bet-hedging (i.e. investing only a fraction of the population to each possible phenotype) offers a good alternative to adapt to this fluctuating environment \cite{kussellPhenotypicDiversityPopulation2005,huftonPhenotypicSwitchingPopulations2018,camachomateuPhenotypicdependentVariabilityEmergence2021,collectiveEconomicPrinciplesCell,nakashimaPopulationDynamicsModels2023}.

In these circumstances our current understanding of evolutionary adaptation suggests that populations tend to evolve toward the strategy which maximizes their mean growth rate
\begin{equation} \label{eq:avg_growth_rate}
    \langle\mu\rangle=\left\langle\frac{1}{N(t)}\frac{dN}{dt}(t)\right\rangle
\end{equation}
as first argued by \cite{gillespieNaturalSelectionVariances1977}. More precisely, \cite{gillespieNaturalSelectionVariances1977} argued that what gets maximized is the geometric mean of the offspring number $\mu'$ (where $N(t+\Delta t)=\mu'\cdot N(t)$) but this is equivalent to maximizing the arithmetic mean of the growth rate $\langle\mu\rangle$ as we have defined it in \autoref{eq:avg_growth_rate}. However, this is not an exact result for finite populations, and when environmental fluctuations are strong and the extinction rate $r_\text{ext}$ (that is, the inverse of the mean extinction time $\tau_\text{ext}$) is non-negligible, as in limited-size populations, we expect to see deviations from this principle.

To date, many studies have explored the probability of fixation of single mutations in fluctuating environments \cite{ashcroftFixationFinitePopulations2014,saakianAlleleFixationProbability2019,kaushikTimeFixationChanging2021}, but not within the context of bet-hedging. Conversely, while other studies have extensively characterized the properties and trade-offs of different bet-hedging strategies \cite{kussellPhenotypicDiversityPopulation2005,huftonPhenotypicSwitchingPopulations2018,dinisParetooptimalTradeoffPhenotypic2022}, they have not looked into what the dynamical process of evolution would actually select.

This gap in the literature is particularly relevant because, as previously reported in \cite{dinisPhaseTransitionsOptimal2020,dinisParetooptimalTradeoffPhenotypic2022} and as we will show in \autoref{sec:no_evolution}, there exist strategies in this context that achieve a noticeable reduction in fluctuations with almost no cost to the average growth rate. The existence of these strategies suggests that the traditional focus on mean growth maximization may be insufficient to predict evolutionary outcomes in finite populations. Exactly which of these bet-hedging strategies is favored by the actual process of natural selection remains an open question, and it is the one we aim to address with this work.

\section{The model: a Markov process} \label{sec:model}

To address here the question outlined in the introduction we introduce the simple agent-based model illustrated in \autoref{fig:model_diagram}. The state of the system consists of only two components: the \emph{environment} and the \emph{population}.

The \emph{environment} is a discrete variable with two possible values (this is easily generalizable to the case of $n_\text{env}$ different environments), $e\in\{1,2\}$, which evolves according to a continuous-time Markov chain with transition rates $\omega_t(e,e')$ between states $e$ and $e'$.

The \emph{population} consists of $N$ agents, each characterized by two traits: a \emph{phenotype} with two possible values, $\phi\in\{A,B\}$, (again, this could easily be generalized to more than two phenotypes) and a probability distribution over these possible phenotypes, $s(\phi)$. We will refer to this distribution as the agent's \emph{phenotypic strategy}. The phenotype will determine the response of the agent to the environment, since the population will follow a stochastic birth-death process in which the birth and death rates ($\omega_b(e,\phi)$ and $\omega_d(e,\phi)$ respectively) depend explicitly on both the current environment and the phenotype of each agent. On the other hand, the phenotypic strategy will regulate the phenotypic composition of the population: upon agent replication, the offspring will exhibit phenotype $A$ with probability $s(A)$ and phenotype $B$ with probability $s(B)$. Since we restrict ourselves to two phenotypes, the strategy is fully specified by a single parameter, $s(A)$, as $s(B)=1-s(A)$. Throughout this work, we will therefore specify only this first component. Thus, although the agents' phenotypes are determined at birth and remain fixed, this mechanism enables phenotypic diversity to emerge and be maintained \emph{across generations} (provided it is favorable to the population).

Additionally, when an agent replicates, its offspring will generally inherit the parent's phenotypic strategy $s$, but with probability $p_\text{mut}$ a mutation will occur, leading to a new random strategy $s_\text{mut}$. This will enable the evolutionary exploration of the space of possible phenotypic strategies, which constitute the heritable traits subject to selection, and therefore play the role of the \emph{genotype}. Thus, for simplicity, mutations in this model are assumed to be abrupt, even though this may not be the most common situation in biological systems. In \autoref{sec:incremental_mutations}, we extend the model by introducing an additional parameter to allow for incremental mutations, and show that this modification leads to qualitatively similar results, with only minor quantitative differences.

The mechanism we have chosen to introduce bet-hedging is different from many previous works \cite{kussellPhenotypicDiversityPopulation2005,huftonPhenotypicSwitchingPopulations2018,dinisParetooptimalTradeoffPhenotypic2022} which used stochastic phenotypic switching. However, this mechanism allows us to describe the strategy with a single variable, allowing for an easier understanding of what strategies are selected while being a very generic and biologically relevant way of modeling bet-hedging.

Finally, as we pointed out in the introduction we expect to see deviations from the already known principles of evolution when the population size does not grow indefinitely. Therefore, we will limit the population size to its initial value, $N_\text{ini}$: whenever replication would cause the population to exceed this value, a randomly chosen agent is simultaneously removed. This mechanism mimics continuous-culture evolution experiments \cite{collectiveEconomicPrinciplesCell}, and can be interpreted as imposing an effective ecological carrying capacity while preserving exponential growth dynamics.

\begin{figure}
    \subfloat[Environment dynamics]{
        \begin{tikzpicture}
            \tikzset{
                environment/.style={draw, rounded corners, minimum size=0.75cm, fill=lightgray!25},
                rate/.style={->, >=stealth, thick},
                node distance=2.75cm
            }
            \node[environment] (e1) {$1$};
            \node[environment, right of=e1] (e2) {$2$};
            \draw[rate, bend left=20] (e1) to node[above] {$\omega_t(1,2)$} (e2);
            \draw[rate, bend left=20] (e2) to node[below] {$\omega_t(2,1)$} (e1);
        \end{tikzpicture}
    }\\
    \subfloat[Birth-death process]{
        \begin{tikzpicture}
            \tikzset{
                agent/.style={draw, rounded corners, minimum height=0.75cm, minimum width=1.25cm, fill=lightgray!25},
                rate/.style={->, >=stealth, thick},
                node distance=2.75cm
            }
            \node[agent] (agent) {$\phi$, $s$};
            \node[agent, right of=agent, yshift=+0.5cm] (death) {$\varnothing$};
            \node[agent, right of=agent, yshift=-0.5cm] (agent_copy) {$\phi$, $s$};
            \node[right=0cm of agent_copy, draw=none] (plus) {$+$};
            \node[agent, right=0cm of plus] (offspring) {$\phi_o$, $s_o$};
            \draw[rate, bend left=10] (agent) to node[sloped, midway, above] {$\omega_d(e,\phi)$} (death);
            \draw[rate, bend right=10] (agent) to node[sloped, midway, below] {$\omega_b(e,\phi)$} (agent_copy);
        \end{tikzpicture}
    }\\
    \subfloat[Phenotype and strategy of the offspring]{
        \begin{tikzpicture}
            \tikzset{
                trait/.style={},
                prob/.style={->},
                node distance=2.25cm
            }
            \node[trait] (phi_o) {$\phi_o=$};
            \node[trait, right of=phi_o, yshift=+0.5cm] (A) {$A$};
            \node[trait, right of=phi_o, yshift=-0.5cm] (B) {$B$};
            \draw[prob] (phi_o) -- node[sloped, midway, above] {$s(A)$} (A);
            \draw[prob] (phi_o) -- node[sloped, midway, below] {$s(B)$} (B);
            \node[trait, right of=phi_o, xshift=1.5cm] (s_o) {$s_o=$};
            \node[trait, right of=s_o, yshift=+0.5cm] (s) {$s$};
            \node[trait, right of=s_o, yshift=-0.5cm] (s_mut) {$s_\text{mut}$};
            \draw[prob] (s_o) -- node[sloped, midway, above] {$1-p_\text{mut}$} (s);
            \draw[prob] (s_o) -- node[sloped, midway, below] {$p_\text{mut}$} (s_mut);
        \end{tikzpicture}
    }
    \caption{Schematic diagram of our model: the \emph{environment} variable follows a simple markov process between two states (a), the population follows a \emph{birth-death process} with rates that depend on the environment and the phenotypes of the agents (b) and finally, the phenotype of the offspring will be determined by the strategy of the parent and the offspring will inherit that same strategy unless a mutation occurs (c).} \label{fig:model_diagram}
\end{figure}

Despite the minimalism of the model, it is analytically intractable. Therefore, we will perform simulations using the Gillespie algorithm \cite{toralStochasticNumericalMethods}. More details about the simulations and the algorithm can be found in \autoref{sec:gillespie_algorithm}. Nonetheless, as we will see in the next sections, the model is simple enough that its behavior can be readily analyzed, and at the same time it contains all the necessary elements to provide insight into the questions that motivate this work.

\section{Behaviour of the model without evolution} \label{sec:no_evolution}

\begin{table}
    \begin{tabular}{c c c}
        \toprule
        symbol             & name             & value                                                        \\
        \midrule
        $\omega_t(e,e')$   & transition rates & $\begin{psmallmatrix}-1.0&+1.0\\+1.0&-1.0\end{psmallmatrix}$ \\
        \addlinespace
        $\omega_b(e,\phi)$ & birth rates      & $\begin{psmallmatrix}1.0&0.2\\0.0&0.0\end{psmallmatrix}$     \\
        \addlinespace
        $\omega_d(e,\phi)$ & death rates      & $\begin{psmallmatrix}0.0&0.0\\1.0&0.1\end{psmallmatrix}$     \\
        \bottomrule
    \end{tabular}
    \caption{Set of parameters used in the simulations of \cref{sec:no_evolution,sec:evol_outcome,sec:heuristic_interpretation,sec:incremental_mutations}, which we will refer to as the \emph{Persister-Resistant} (PR) parameter set.} \label{tab:PR_parameters}
\end{table}

\begin{figure}
    \includegraphics{../sims/asymmetric/plots/strat_phe_0_i/rates.pdf}
    \caption{Extinction rate $r_\text{ext}$ versus mean growth rate $\langle\mu\rangle$. Results are shown for \texttt{fixed}, \texttt{evol} and \texttt{evol(r)} simulations, across a range of initial strategies from $s(A)=0$ (lower end of the curve) to $s(A)=1$ (upper end). These values define a \emph{Pareto front} capturing the trade-off between growth and survival.} \label{fig:pareto_front}
\end{figure}

We begin by solving the model without evolution, i.e. with $p_\text{mut}=0$, setting the initial population size to $N_\text{ini}=100$ and assuming that all agents initially share the same strategy, $s_\text{ini}$. This kind of simulations (which we will call \texttt{fixed}) allow us to study the properties of individual strategies in isolation and, in turn, better understand the evolutionary outcomes that follow.

To proceed, we must first choose a specific set of model parameters. We start with the \emph{Persister-Resistant} (PR) parameter set shown in \autoref{tab:PR_parameters} (a completely different set of parameters is considered in \autoref{sec:symmetric_parameters}, which leads to similar results). These parameters represent a common biological scenario in which there is a ``persister'' phenotype and a ``resistant'' one \cite{huftonPhenotypicSwitchingPopulations2018}. More precisely the meaning of this parameters is the following: both environments are equally likely and there is a transition rate between them of $1.0$; environment $1$ represents a ``favorable'' environment in which phenotype $A$ (the ``resistant'' one) replicates with a rate of $1.0$ while phenotype $B$ (the ``persister'' one) does it with a rate of $0.2$, and neither of them die; environment $2$ instead represents ``harsh'' conditions in which phenotype $A$ dies with a rate of $1.0$ and phenotype $B$ with a rate of $0.1$, and neither of them are able to duplicate.

\autoref{fig:pareto_front} shows the extinction rate $r_\text{ext}$ versus mean growth rate $\langle\mu\rangle$ in simulations with the PR parameter set. Here, we restrict our attention to the \texttt{fixed} simulations (represented by the blue curve) with initial strategies ranging from $s(A)=0$ (lower end of the curve) to $s(A)=1$ (upper end). The first important observation is that the maximum of $\langle\mu\rangle$ is achieved neither at $s(A)=0$ nor at $s(A)=1$, but rather at an intermediate strategy, indicating that in scenarios such as this one phenotypic diversity (that is, maintaining a mixture of both phenotypes) is the ``optimal'' strategy, a result already reported multiple times \cite{kussellPhenotypicDiversityPopulation2005,xueBetHedgingDemographic2017,dinisPhaseTransitionsOptimal2020,dinisParetooptimalTradeoffPhenotypic2022}. As we will see, this is precisely the strategy reached by evolution at large $N_\text{ini}$. This behavior does not arise in all cases, but only when environmental fluctuations make each phenotype individually less favorable in the long term than their combination. Secondly, we observe that the growth-maximizing strategy differs from the extinction-minimizing one (again, this holds true for other parameter sets, as discussed in \autoref{sec:symmetric_parameters}). In particular, $\langle\mu\rangle$ is maximized at $s(A)\sim0.8$, while $r_\text{ext}$ is minimized at $s(A)\sim0.0$. As we will show in \autoref{sec:heuristic_interpretation}, this trade-off (and more precisely the different scaling of these two quantities with the population limit) plays a key role in shaping the evolutionary outcome. Note that all figures throughout this work include error bars or error bands; however, in many cases, they are smaller than the point symbols and are therefore barely visible.

It is important to clarify that many of the quantities analyzed throughout this work (like $\langle\mu\rangle$ in \autoref{fig:pareto_front}) represent time averages and their estimation requires good statistics. However, since populations eventually undergo extinction under these stochastic dynamics (often on timescales too short to gather sufficient data from a single run), we concatenate multiple independent realizations into a single, longer trajectory. All time averages and distributions are subsequently computed across these extended, joint trajectories.\footnote{\textcolor{red}{Creo que sería interesante explorar y considerar otros tipos de promedio, como por ejemplo promedios calculados a un tiempo fijo concreto $t$ medido desde el inicio de la simulación.}}

Rather than exploring in detail how these results depend on the specific choice of parameters (a topic extensively studied in previous work \cite{kussellPhenotypicDiversityPopulation2005,dinisParetooptimalTradeoffPhenotypic2022}) the next section focuses on the evolutionary dynamics and the strategies ultimately selected under these conditions.

\section{Evolutionary dynamics} \label{sec:evol_outcome}

We now enable evolution by setting $p_\text{mut}=0.001$ (while keeping $N_\text{ini}=100$), allowing strategies to mutate and compete. We consider two types of simulations: in \texttt{evol} simulations, the population is initially composed of agents with the same strategy $s(A)_\text{ini}$, while in \texttt{evol(r)} simulations, the initial strategy of each agent is completely random.

\begin{figure}
    \includegraphics{../sims/asymmetric/plots/strat_phe_0_i/avg_avg_strat_phe_0.pdf}
    \caption{Time average of the population-averaged strategy, $\overline{s}(A)$, as a function of the initial strategy $s_\text{ini}$ for the PR parameter set. Results are shown for \texttt{fixed}, \texttt{evol} and \texttt{evol(r)} simulations (represented as a horizontal line, as they lack a single defined initial strategy). The dashed and dotted lines indicate the strategies that maximize mean growth and minimize extinction, respectively, with the evolved strategy lying between them.} \label{fig:evol_results}
\end{figure}

To characterize the outcome of the evolutionary process, we compute the population-averaged strategy $\overline{s}(A)$ and analyze its time average, denoted as $\langle\overline{s}\rangle(A)$. In \autoref{fig:evol_results}, we show $\langle\overline{s}\rangle(A)$ as a function of the initial strategy $s_\text{ini}$. As expected, for the \texttt{fixed} simulations, the initial strategy is preserved, producing a diagonal line. In contrast, for the evolutionary simulations, the system consistently converges toward a well-defined stationary value, exhibiting only a weak dependence on the initial condition.

A key observation is that the instantaneous dispersion of strategies $\sigma_s$ (defined as the standard deviation of strategies present in the population at a given time) remains small. This dispersion is represented in \autoref{fig:evol_results} as an error band, which is barely visible. At any given time, the population is therefore effectively characterized by a single dominant strategy. Over longer times, however, mutation-driven exploration causes this dominant strategy to drift through the space of possible strategies as can be seen in \autoref{fig:time_series}. More importantly, we observe that the average evolved strategy lies between the growth-maximizing and extinction-minimizing strategies identified in the fixed-strategy simulations. As we will see in \autoref{sec:symmetric_parameters}, this behavior is not specific to this parameter set. This shows that evolution in finite populations does not optimize instantaneous growth alone, but instead selects strategies that represent a compromise between growth and extinction avoidance.

\begin{figure}
    \includegraphics{../sims/asymmetric/plots/time_series/avg_strat_phe_0.pdf}
    \caption{Time series of the population-averaged strategy over multiple concatenated simulations (extinction events are marked by vertical dotted lines). The error band displays the strategy standard deviation within the population at each instant, $\sigma_s$, which remains remarkably narrow. Notably, before extinction is reached, advantageous mutations frequently emerge and fix, abruptly shifting the average strategy and driving a broad exploration of the strategy space.} \label{fig:time_series}
\end{figure}

Finally, in \autoref{fig:dependence_on_N_ini_and_p_mut} we explore how the evolutionary outcome (namely, the average strategy $\langle\overline{s}\rangle(A)$) depends on the population size $N_\text{ini}$ and the mutation probability $p_\text{mut}$. We first observe a strong dependence on $N_\text{ini}$. As anticipated, when $N_\text{ini}$ increases, $\langle\overline{s}\rangle(A)$ approaches the growth-maximizing value identified in \autoref{sec:no_evolution}. Only for small populations does the evolutionary outcome favor safer strategies with lower extinction rates. In contrast, the dependence on the mutation probability is relatively weak, as $p_\text{mut}$ mainly controls the exploration of the strategy space (in \autoref{sec:incremental_mutations} we further examine the effect of changing the characteristics of mutations). For very large values of $p_\text{mut}$, the strategy standard deviation $\sigma_s$ increases, indicating that the population no longer shares a common strategy at a given time. In fact, in the limit $p_\text{mut} \to 1$, the population behaves as if $s(A) = 0.5$, since strategies change randomly with every replication. Conversely, in the limit $p_\text{mut} \to 0$, mutations are absent, and the outcome is determined solely by selection among the strategies present in the initial condition.

\begin{figure}
    \includegraphics{../sims/asymmetric/plots/n_agents_i/avg_avg_strat_phe_0.pdf}
    \includegraphics{../sims/asymmetric/plots/prob_mut/avg_avg_strat_phe_0.pdf}
    \caption{Dependence of the average evolved strategy on the initial population size $N_\text{ini}$ (top) and the mutation probability $p_\text{mut}$ (bottom). Increasing $N_\text{ini}$ drives the system toward the growth-maximizing strategy, while smaller populations favor safer strategies with lower extinction rates. The dashed and dotted lines indicate the growth-maximizing and extinction-minimizing strategies, respectively, and the error band displays the dispersion of strategies $\sigma_s$.} \label{fig:dependence_on_N_ini_and_p_mut}
\end{figure}

\section{Heuristic interpretation of the evolutionary outcome} \label{sec:heuristic_interpretation}

To better understand the outcome of the evolutionary process, it is useful to examine in more detail the trade-off between the mean growth rate and the extinction rate. As discussed in \autoref{sec:no_evolution}, these two quantities cannot be optimized simultaneously. If we plot, as in \autoref{fig:pareto_front}, the values of $r_\text{ext}$ against $\langle\mu\rangle$ for \texttt{fixed} simulations with different intial strategies, we obtain a curve known as a Pareto front \cite{dinisPhaseTransitionsOptimal2020}. This curve represents the set of optimal trade-offs: any increase in growth necessarily comes at the cost of a higher extinction risk, and vice versa. In more general settings with more than two phenotypes, strategies would span a higher-dimensional space, and the Pareto front would correspond to the boundary (also known as the efficient frontier) of the accessible region.

In the same plot we also display the outcome obtained when evolution is enabled (both for \texttt{evol} and \texttt{evol(r)} simulations), allowing us to compare directly the resulting growth and extinction rates rather than the underlying strategy. We find that the system converges (since the values of $r_\text{ext}$ and $\langle\mu\rangle$ concentrate), with only a moderate dependence on the initial conditions, toward a point with a lower growth rate than the growth-maximizing strategy, but also with a reduced extinction rate.

Importantly, this point does not lie exactly on the Pareto front defined by fixed strategies, indicating that the evolutionary outcome cannot be associated with any single fixed strategy. Instead, both the mean growth rate and the extinction rate are well approximated by averages over the distribution of strategies:
\begin{equation} \label{eq:evol_averages}
    \langle\mu\rangle\approx\int\langle\mu\rangle(\overline{s})p(\overline{s})d\overline{s}, \quad r_\text{ext}\approx\int r_\text{ext}(\overline{s})p(\overline{s})d\overline{s}
\end{equation}
where $\langle\mu\rangle(s)$ and $r_\text{ext}(s)$ represent the mean growth rate and extinction rate that a \emph{fixed} strategy $s$ has, while $p(\overline{s})$ is the distribution of population-averaged strategies in evolutive simulations. In \autoref{sec:supplementary_results} we verify that this holds across a wide range of parameters.

This behavior can be understood from the dynamical nature of the evolutionary process: at any given time, the population is effectively characterized by a single dominant strategy, but this strategy slowly changes due to mutations and reinitializations following extinction events. Over long times, the system can therefore be interpreted as a concatenation of fixed-strategy dynamics with different values of $\overline{s}$.

\begin{figure}
    \includegraphics{../sims/asymmetric/plots/fixed/avg_growth_rate.pdf}
    \includegraphics{../sims/asymmetric/plots/fixed/extinct_rate_scaling.pdf}
    \caption{Fixed-strategy results used to construct the heuristic approximation. The mean growth rate (top) depends only weakly on $N_\text{ini}$, whereas the extinction rate (bottom) follows an approximate power-law scaling. Dashed lines indicate spline interpolations of the simulation data.} \label{fig:fixed_results}
\end{figure}

\begin{figure}
    \includegraphics{../sims/asymmetric/plots/fixed/exp_avg_avg_strat_phe_0.pdf}
    \includegraphics{../sims/asymmetric/plots/fixed/avg_avg_strat_phe_0.pdf}
    \caption{Expected average strategy from the heuristic distribution $\langle\overline{s}(A)\rangle^\ast$ (top) and observed average strategy from \texttt{evol(r)} simulations (bottom) for different values of $N_\text{ini}$ and $p_\text{mut}$. The heuristic approximation captures the qualitative dependence of the evolutionary outcome on $N_\text{ini}$ and $p_\text{mut}$, albeit with some quantitative discrepancies.} \label{fig:heuristic_approx}
\end{figure}

The results above motivate the search for a heuristic approximation of the distribution $p(\overline{s})$, constructed solely from properties of fixed strategies, that could allow us to predict the evolutionary outcome $\langle\overline{s}\rangle(A)$ as a function of $N_\text{ini}$ and $p_\text{mut}$. To construct this approximation, we will begin by assuming that the probability of a given strategy being selected during an initial competitive phase following reinitialization is proportional to $\exp(N_\text{ini}\cdot\langle\mu\rangle(\overline{s}))$. This is because differences in the average growth rate will be exponentially amplified over the duration of this competitive phase and the time required to resolve this competition will scale with the number of competing strategies (as a first approximation, we will assume this duration to be directly proportional to $N_\text{ini}$, which justifies its presence in the exponent). However, to obtain the long-term distribution of strategies, this selection probability must be weighted by the expected survival time before the population goes extinct or is replaced by a mutant. This leads us to divide the exponential term by the sum of two distinct rates: the population extinction rate, $r_\text{ext}(N_\text{ini},\overline{s})$, and the rate of mutation-driven replacement. While an exact derivation of the latter is generally prohibitive, it acts as a minor correction where an order-of-magnitude estimate is sufficient. This replacement rate can be estimated as the total mutation emergence rate, $N_\text{ini}p_\text{mut}\langle\mu_b\rangle(\overline{s})$ (where $\langle\mu_b\rangle$ is the average birth rate, excluding deaths), multiplied by the fixation probability of the new mutant. Although calculating this probability exactly is highly challenging \cite{ashcroftFixationFinitePopulations2014,saakianAlleleFixationProbability2019,kaushikTimeFixationChanging2021}, it is of the order of $1/N_\text{ini}$ and thus, $N_\text{ini}$ cancels out. Together, these terms yield the following heuristic form for the distribution of strategies:
\begin{equation} \label{eq:heuristic_distribution}
    p(\overline{s})\propto\frac{\exp(N_\text{ini}\cdot\langle\mu\rangle(\overline{s}))}{r_\text{ext}(N_\text{ini},\overline{s})+p_\text{mut}\cdot\langle\mu_b\rangle(\overline{s})}.
\end{equation}
To evaluate $\langle\mu\rangle$, $\langle\mu_b\rangle$ and $r_\text{ext}$ for arbitrary values of $\overline{s}$ and $N_\text{ini}$ we use spline interpolations of the numerical data obtained from \texttt{fixed} simulations.

To understand how this distribution will depend on the model parameters, it is useful to examine how $\langle\mu\rangle(s)$ and $r_\text{ext}(N_\text{ini},s)$ depend on the population size, as shown in \autoref{fig:fixed_results}. We find that the mean growth rate depends only weakly on $N_\text{ini}$, except for very small populations, making it primarily a function of $s$. In contrast, the extinction rate exhibits a strong dependence, decreasing approximately as a power law, $r_\text{ext}\sim N_\text{ini}^{-\alpha}$. This scaling arises because extinctions are primarily driven by unfavorable environmental episodes. In an unfavorable environment we will have (on average) $N(t)\sim N_\text{ini}e^{-\gamma t}$, meaning that reaching extinction requires an adverse episode of duration $T\sim\ln(N_\text{ini})/\gamma$. Since environmental durations are exponentially distributed, the probability of encountering such a long episode scales as $e^{-\lambda T}$, which leads to the power-law dependence $r_\text{ext}\sim N_\text{ini}^{-\lambda/\gamma}$. Consequently, in \autoref{eq:heuristic_distribution}, the growth contribution enters exponentially through the numerator, whereas extinction risk enters algebraically in the denominator. This scaling mismatch is the main reason behind the transition: extinction risk dominates at small population sizes, favoring safer strategies, while growth maximization inevitably prevails in the large-$N_\text{ini}$ limit.

Finally, \autoref{fig:heuristic_approx} compares the expected average strategy calculated from the heuristic distribution (top) with the actual average observed in evolutionary simulations (bottom). While a detailed comparison of the full probability distributions is deferred to \autoref{sec:supplementary_results}, focusing on their mean values reveals good qualitative agreement across a wide range of $N_\text{ini}$ and $p_\text{mut}$ values. Crucially, the heuristic successfully captures the transition from growth-maximizing to safer strategies as $N_\text{ini}$ decreases, as well as how this transition is modulated by $p_\text{mut}$.

\section{Incremental mutations} \label{sec:incremental_mutations}

We also consider a modified mutation mechanism in which changes to the strategy are incremental rather than completely random. Specifically, when a mutation occurs, the offspring strategy is obtained by adding Gaussian noise with standard deviation $\sigma_\text{mut}$ to the parent's strategy, followed by normalization.

For $\sigma_\text{mut}=0.1$ and $p_\text{mut}=0.01$, the results are shown in \autoref{fig:incremental_mutations}. The qualitative behavior remains unchanged: evolution still leads to intermediate strategies, and the dependence on population size is preserved.

However, we observe a slightly stronger dependence on initial conditions, reflecting the slower exploration of strategy space induced by incremental mutations.

\begin{figure}
    \includegraphics{../sims/incremental/plots/strat_phe_0_i/avg_avg_strat_phe_0.pdf}
    \includegraphics{../sims/incremental/plots/n_agents_i/avg_avg_strat_phe_0.pdf}
    \caption{Results for incremental mutations. Left: dependence of the evolved strategy on initial conditions. Right: dependence on population size. The qualitative behavior matches that of the main model.} \label{fig:incremental_mutations}
\end{figure}

\section{Different model parameters} \label{sec:symmetric_parameters}

We also analyzed a qualitatively different parameter regime defined by:
\begin{itemize}
    \item $\omega_b(e,\phi)=\begin{psmallmatrix}1.2&0.0\\0.0&0.8\end{psmallmatrix}$
    \item $\omega_d(e,\phi)=\begin{psmallmatrix}0.0&1.0\\1.0&0.0\end{psmallmatrix}$
\end{itemize}

In this case, each phenotype is specialized for one environment, representing a scenario fundamentally different from the asymmetric case discussed in the main text.

As shown in \autoref{fig:symmetric_parameters}, the trade-off between growth and extinction persists, although the Pareto front is much narrower: the strategies that maximize growth and minimize extinction are closer to each other.

Despite these differences, the main conclusions remain unchanged. Evolution leads to a well-defined strategy largely independent of initial conditions, lying between the growth-maximizing and extinction-minimizing values. As seen in \autoref{fig:symmetric_parameters}, the population size $N_\text{ini}$ again controls the outcome: large populations favor growth-maximizing strategies, while small populations favor safer ones.

Finally, our analysis confirmed that the heuristic approximation still captures the qualitative transition between regimes. However, quantitative discrepancies are more pronounced in this symmetric case: in particular, the convergence toward the growth-maximizing strategy occurs at significantly larger $N_\text{ini}$ than predicted.

\begin{figure}
    \includegraphics{../sims/symmetric/plots/strat_phe_0_i/avg_avg_strat_phe_0.pdf}
    \includegraphics{../sims/symmetric/plots/n_agents_i/avg_avg_strat_phe_0.pdf}
    \caption{Average evolved strategy as a function of the initial strategy (left) and dependence of the evolved strategy on population size (right) in the symmetric parameter case. The dashed line indicates the strategy that maximizes mean growth, while the dotted line marks the one that minimizes extinction. The same qualitative trends as in the asymmetric case are observed.} \label{fig:symmetric_parameters}
\end{figure}

\section{Discussion and outlook} \label{sec:discussion}

Our results provide a clear and minimal illustration of how evolutionary outcomes in finite populations cannot be inferred solely from the average instantaneous growth rate. The long-term dominant strategy is shaped not only by the mean behavior of the population but also by fluctuations, particularly those that can lead to extinction. In fluctuating environments with a population size limit, extinction is inevitable over long times and thus acts as a genuine selective pressure.

In the large-$N_\text{ini}$ limit, extinction risk becomes negligible, and evolution correctly selects the strategy that maximizes $\langle\mu\rangle$. Importantly, this growth-maximizing strategy is neither $s(A)=0$ nor $s(A)=1$, highlighting that maintaining phenotypic diversity can be advantageous in fluctuating environments. For smaller populations, however, extinction risk reshapes the adaptive landscape, biasing evolution toward safer, lower-growth strategies. This crossover demonstrates how demographic constraints fundamentally alter evolutionary optimization.

The heuristic approximation introduced here captures the qualitative competition between growth and extinction, clarifying the scaling mechanisms underlying the transition between growth-maximizing and safe strategies. Quantitative discrepancies remain, emphasizing that evolutionary dynamics cannot be fully reduced to fixed-strategy observables. Extending this framework toward a more systematic coarse-grained theory, possibly linked to quasi-stationary distributions, represents an important avenue for future work.

More broadly, this study provides a flexible framework to explore evolution in fluctuating environments. The model can be generalized to more environments or phenotypes, and to strategies that respond to environmental cues (sensing), enabling the systematic study of bet-hedging, plasticity, and stochastic selection in complex settings.

\begin{acknowledgments}
    We wish to acknowledge financial support through grant PID2023-147067NB-I00, funded by the Agencia Española de Investigación (AEI, Spain) and the Ministerio de Ciencia, Innovación y Universidades (MICIU, Spain) MICIU/AEI/10.13039/501100011033, and by the European Regional Development Fund (ERDF, European Union).
\end{acknowledgments}

\bibliographystyle{apsrev4-2}
\bibliography{mutare.bib}

\appendix

\section{The Gillespie algorithm and time series analysis} \label{sec:gillespie_algorithm}

To simulate the stochastic dynamics of the model, we use the Gillespie algorithm, a standard method for generating statistically exact trajectories of continuous-time Markov processes.

At any given time, the state of the system is denoted by $\mathbf{X}$ and all possible events (birth, death, environmental transitions) are indexed by $j=1,\dots,M$, each occurring with rate $\omega_j(\mathbf{X})$. The total rate is therefore
\begin{equation}
    \Omega(\mathbf{X})=\sum_{j=1}^M\omega_j(\mathbf{X}).
\end{equation}

The algorithm proceeds iteratively as follows. First, the time $\tau$ to the next event is drawn from an exponential distribution:
\begin{equation}
    \tau=\frac{1}{\Omega(\mathbf{X})}\ln\frac{1}{r_1},\quad r_1\sim U(0,1).
\end{equation}
Then, the specific event $\mu$ that occurs is selected according to:
\begin{equation}
    \sum_{j=1}^{\mu-1}\omega_j(\mathbf{X})<r_2\Omega(\mathbf{X})\le\sum_{j=1}^{\mu}\omega_j(\mathbf{X}),\quad r_2\sim U(0,1).
\end{equation}
Finally, the system is updated as
\begin{equation}
    \mathbf{X}(t+\tau)=\mathbf{X}(t)+\boldsymbol{\nu}_\mu,
\end{equation}
where $\boldsymbol{\nu}_\mu$ is the state-change vector associated with event $\mu$.

An important consequence of this algorithm is that observations are naturally sampled at irregular time intervals $\Delta t_k=\tau_k$. Therefore, averages of observables must be computed as time-weighted averages rather than simple arithmetic means.

In addition to the mean population growth rate $\langle\mu\rangle$, we also measure fluctuations of growth along stochastic trajectories. A naive definition based on the instantaneous rate $\Delta N/(N\,\Delta t)$ leads to ill-defined fluctuations due to arbitrarily small waiting times. Instead, we define fluctuations through the integrated growth process.

Denoting by
\begin{equation}
    \mu_k=\frac{\Delta N_k}{N_k\,\Delta t_k}
\end{equation}
the growth rate measured over a Gillespie step $k$, we compute
\begin{equation}
    \sigma_\mu^2=\frac{1}{\sum_k \Delta t_k}\sum_k\left(\mu_k-\langle\mu\rangle\right)^2\Delta t_k^2.
\end{equation}
This corresponds to the variance growth rate of the cumulative per-capita growth,
\begin{equation}
    \log\frac{N(t)}{N(0)}=\sum_k \mu_k\,\Delta t_k,
\end{equation}
and is therefore a well-defined quantity in continuous-time jump processes. In this sense, $\sigma_\mu^2$ should be interpreted as the diffusion coefficient of log-population growth rather than as the variance of an instantaneous rate.

\section{Supplementary Results} \label{sec:supplementary_results}

\begin{figure*}
    \includegraphics{../sims/asymmetric/plots/strat_phe_0_i/avg_birth_rate.pdf}
    \includegraphics{../sims/asymmetric/plots/strat_phe_0_i/avg_dist_phe_0.pdf}
    \caption{Average birth rate (left) and phenotype distribution (right) as functions of the strategy. The observed phenotype composition lies between the theoretical equilibrium predictions for each environment.} \label{fig:supplementary_observables_1}
\end{figure*}

\begin{figure*}
    \includegraphics{../sims/asymmetric/plots/strat_phe_0_i/dist_n_agents.pdf}
    \includegraphics{../sims/asymmetric/plots/strat_phe_0_i/std_dev_growth_rate.pdf}
    \caption{Distribution of population sizes (left) and standard deviation of growth rate (right) as functions of the strategy. Increased weight at low population sizes correlates with higher extinction risk.} \label{fig:supplementary_observables_2}
\end{figure*}

\begin{figure*}
    \includegraphics{../sims/asymmetric/plots/time_series/avg_strat_phe_0.pdf}
    \includegraphics{../sims/asymmetric/plots/prob_mut/avg_growth_rate.pdf}
    \caption{Left: time series of the population-averaged strategy showing dominance of a single strategy at each time. Right: comparison illustrating the validity of Eq.~\eqref{eq:evol_averages}.} \label{fig:evidence_of_common_strategy}
\end{figure*}

\begin{figure*}
    \includegraphics{../sims/asymmetric/plots/fixed/exp_dist_avg_strat_phe_0.pdf}
    \includegraphics{../sims/asymmetric/plots/fixed/dist_avg_strat_phe_0.pdf}
    \caption{Comparison between heuristic (left) and simulated (right) distributions of strategies for different parameters. Qualitative agreement is observed, but differences remain.} \label{fig:strategy_distributions}
\end{figure*}

Here we present additional observables that complement the results shown in the main text.

Among them, we consider the average birth rate, which increases monotonically with the fraction of the resistant phenotype, unlike the total growth rate. We also analyze the distribution of phenotypes in the population, which differs from the strategy due to phenotype-dependent lifetimes. As shown in \autoref{fig:supplementary_observables_1}, this distribution lies between the equilibrium predictions corresponding to each environment.

These equilibrium predictions can be obtained from the dominant eigenvector of
\begin{equation}
    \begin{pmatrix}s(0)\omega_b(e,0)-\omega_d(e,0)&s(0)\omega_b(e,1)\\s(1)\omega_b(e,0)&s(1)\omega_b(e,1)-\omega_d(e,1)\end{pmatrix},
\end{equation}
which describes the deterministic dynamics in a fixed environment.

We also examine the distribution of population sizes and the fluctuations of the growth rate (see \autoref{fig:supplementary_observables_2}). The population size distribution is typically bimodal, and the prominence of the low-population peak correlates with extinction risk.

Time series data (\autoref{fig:evidence_of_common_strategy}) confirm that, at low mutation rates, the population is effectively characterized by a single dominant strategy at each instant. The apparent spread of strategies arises from temporal averaging rather than coexistence.

Finally, in \autoref{fig:strategy_distributions}, we compare the full distributions of strategies predicted by the heuristic and observed in simulations. While they share qualitative features, they do not coincide exactly.

\end{document}
